\section{Practical Integration and Research Agenda}
\label{sec:agenda}
\subsection{Choose a Workflow from the Testing Objective}
Workflow selection starts with the behavioral decision, available evidence, and missing artifact. Regression augmentation can use existing tests and coverage, as in TestGen-LLM~\cite{U07}. Conformance needs requirement-backed judgments; fuzzing needs valid generators and failure attribution; interactive code generation can use user-confirmed examples, as in TiCoder~\cite{T11}. Assertion provenance connects these choices to review and maintenance.

\subsection{Industrial Adoption and Human Revision}
Industrial evidence includes different acceptance decisions. TestGen-LLM studies improvement of existing tests through operational filtering~\cite{U07}; the acceptance-testing case study evaluates both generated scenarios and their conversion into executable scripts~\cite{R03}. RAGTAG and system-level test-specification generation additionally involve domain interpretation and assessment~\cite{R01,R02}. ASTER adds dependency handling and mocking, while AToMIC connects industrial tickets to mobile acceptance scripts~\cite{U12,R06}. The OSS-Fuzz-Gen report adds failure-validation work to the adoption burden~\cite{F09}. A scenario judged useful by an engineer and a script accepted into a regression workflow are distinct outcomes.

A deployment evaluation should retain these stages: generated, reviewed, revised, executed, and adopted. Human effort includes interpreting requirements, supplying missing context, correcting assertions, and maintaining the accepted test. Reporting the time and rejection reasons at each stage would connect the artifact's technical quality to its actual workflow value. The reviewed examples motivate these dimensions; this synthesis does not pool their heterogeneous adoption measures.

\subsection{Oracle Provenance and Independent Validation}
Linking an assertion to a requirement clause, validated example, formal property, or regression observation could guide acceptance and repair. The research question is whether checked provenance improves oracle quality and maintenance. A matched-budget comparison should measure unsupported assertions, rejection of legal implementations, independent detection, and review time. Plausible explanations without verified evidence links do not establish that benefit.

\subsection{Feedback That Preserves Valid Tests}
Repair should recover invalid tests while preserving valid judgments. Branching from the same initial artifact allows separate measurement of recovery and regression, followed by independent suite evaluation. Controls can compare setup-only repair, unrestricted repair, no semantic repair, and trusted-reference checking. The information and budget must match; improved passing alone cannot determine whether assertions were corrected or weakened.

\subsection{Generalization in Learned and Adversarial Testers}
Adversarial curricula offer difficult errors~\cite{T02,T03,T07}, but transferable capability requires a frozen tester evaluated on new opponents, repositories, and specification styles. Suite rewards also obscure which tests help or mislead. Comparing aggregate rewards with verified test-level feedback can test the value of finer credit assignment, while retaining validity constraints and independent detection.

\subsection{Repository Scale and Lifecycle Evaluation}
TestGenEval, CITYWALK, and service-level testing provide repository-oriented starting points~\cite{B01,U08,A02}. Lifecycle evaluation should additionally label whether a change preserves semantics or deliberately changes requirements. Evolution and database-flakiness studies motivate retaining prior tests and measuring execution stability, maintenance effort, and unintended-regression detection~\cite{E02,E03}. Version-aware retrieval should be compared with unrestricted context reuse, since existing tests can preserve brittle setup or obsolete obligations.

\subsection{Reproducible Evidence and Efficient Allocation}
Reproducibility packages should expose the path from request to retained suite. Relevant records include model and prompt versions, context selection, generated candidates, tool feedback, repairs, rejection reasons, and final evaluation inputs. This enables analysis of why one method succeeds and whether its filters hide a large failure burden. For training methods, data construction and reward implementation are equally important.

Budget allocation is another cross-cutting research question. A fixed budget can be spent on generating more candidates, checking oracles, executing additional faults, or resolving uncertain requirements. Dynamic scaling in CodeRM~\cite{T06} and reusable fuzzing infrastructure~\cite{F05,F06,F07} illustrate different responses to this allocation problem. A broader study could compare these investments under a common deployment objective, such as maximizing independent detection while bounding invalid-suite rate and latency.
